\documentclass[12pt,a4paper]{report}

\setcounter{section}{0}
\renewcommand{\thesection}{\arabic{section}}
\renewcommand{\thesubsection}{\thesection.\arabic{subsection}}
\renewcommand{\thesubsubsection}{\thesubsection.\arabic{subsubsection}}

\usepackage[utf8]{inputenc}
\usepackage{graphicx}
\usepackage{geometry}
\usepackage{setspace}
\usepackage{mathptmx}
\usepackage{caption}
\usepackage{float}
\usepackage{hanging}
\usepackage{hyperref}
\usepackage{booktabs}
\usepackage{longtable}
\usepackage{array}
\usepackage{enumitem}
\usepackage{fancyhdr}
\usepackage{titlesec}
\usepackage{parskip}

% Page geometry (guidelines: Top 1.5", Bottom 1.0", Left 2.0", Right 1.0")
\setlength{\headheight}{14.5pt}
\geometry{
  a4paper,
  top=1.5in,
  bottom=1.0in,
  left=1.0in,
  right=1.0in
}

\onehalfspacing
\setlength{\parskip}{6pt}

% Section heading fonts (Arial-like via sans-serif) 
\titleformat{\section}{\sffamily\bfseries\fontsize{16}{20}\selectfont}{\thesection.}{0.5em}{}
\titleformat{\subsection}{\sffamily\bfseries\fontsize{14}{18}\selectfont}{\thesubsection}{0.5em}{}
\titleformat{\subsubsection}{\sffamily\bfseries\itshape\fontsize{13}{17}\selectfont}{\thesubsubsection}{0.5em}{}

\pagestyle{fancy}
\fancyhf{}
\fancyhead[RO]{Mahir on Call: Advanced Elucidation Support in Local Language}
\fancyhead[RE]{\leftmark}
\fancyfoot[R]{\thepage}
\renewcommand{\headrulewidth}{0.4pt}

\hypersetup{
  colorlinks=true,
  linkcolor=black,
  urlcolor=blue,
  citecolor=black
}

\begin{document}

\begin{titlepage}
  \thispagestyle{empty}
  \centering
  {\setstretch{1.1}
    \fontsize{26}{32}\selectfont\bfseries Mahir on Call\\[0.4cm]
    \fontsize{20}{26}\selectfont\bfseries Advanced Elucidation Support in Local Language
  }\\[2.0cm]

  {\large Session 2022--2026}\\[0.5cm]
  {\large \textbf{Supervised by}}\\
  {\large Muhammad Umer Haroon}\\[0.4cm]
  {\large \textbf{Co-supervised by}}\\
  {\large Shahzeb Khan}\\[2.0cm]

  \vfill
  \includegraphics[width=0.28\textwidth]{universityLogo.png}\\[1.2cm]

  {\large Department of Computer Science}\\[0.4cm]
  {\Large National University of Computer and Emerging Sciences}\\[0.3cm]
  {\Large Peshawar, Pakistan}
\end{titlepage}

\pagenumbering{roman}
\setcounter{page}{1}

% Student's Declaration
\newpage
\thispagestyle{plain}
\begin{center}
  {\Large \textbf{Student's Declaration}}\\[1cm]
\end{center}

\noindent We declare that this project titled \textit{Mahir on Call: Advanced
Elucidation Support in Local Language}, submitted as requirement for the award
of degree of Bachelors in Computer Science, does not contain any material
previously submitted for a degree in any university; and that to the best of
our knowledge, it does not contain any materials previously published or
written by another person except where due reference is made in the text.\\[0.5cm]

\noindent We understand that the management of the Department of Computer
Science, National University of Computer and Emerging Sciences, has a
zero-tolerance policy towards plagiarism. Therefore, we, as authors of the
above-mentioned thesis, solemnly declare that no portion of our thesis has
been plagiarized, and any material used in the thesis from other sources is
properly referenced.\\[0.5cm]

\noindent We further understand that if we are found guilty of any form of
plagiarism in the thesis work even after graduation, the University reserves
the right to revoke our BS degree.\\[1.2cm]

\begin{tabular}{p{7cm} p{6cm}}
  Uzair Ahmad      & Signature: \hrulefill \\[1cm]
  Arsalan Mateen   & Signature: \hrulefill \\[1cm]
  Muhammad Sohaib  & Signature: \hrulefill \\
\end{tabular}

\vfill
\begin{flushright}
  \begin{tabular}{l}
    \hrulefill \\
    Verified by Plagiarism Cell Officer \\
    Dated: \underline{\hspace{4cm}} \\
  \end{tabular}
\end{flushright}

% Certificate of Approval
\newpage
\thispagestyle{plain}
\begin{center}
  {\Large \textbf{Certificate of Approval}}\\[0.8cm]
  \includegraphics[width=0.25\textwidth]{universityLogo.png}\\[0.8cm]
\end{center}

\noindent The Department of Computer Science, National University of Computer
and Emerging Sciences, accepts this thesis titled \textit{Mahir on Call ---
Advanced Elucidation Support in Local Language}, submitted by Uzair Ahmad,
Arsalan Mateen, and Muhammad Sohaib, in its current form, as fulfilling the
dissertation requirements for the award of Bachelors Degree in Computer
Science.\\[1.5cm]

\begin{center}
  \makebox[0.55\linewidth]{\hrulefill}\\
  Prof. Muhammad Umer Haroon\\
  FYP Supervisor\\
  National University of Computer and Emerging Sciences, Peshawar\\[1.5cm]

  \makebox[0.55\linewidth]{\hrulefill}\\
  Riaz Nawab\\
  FYP Coordinator\\
  National University of Computer and Emerging Sciences, Peshawar\\[1.5cm]

  \makebox[0.55\linewidth]{\hrulefill}\\
  Dr. Qasim Jaan\\
  HoD, Department of Computer Science\\
  National University of Computer and Emerging Sciences, Peshawar
\end{center}

% Acknowledgement 
\newpage
\thispagestyle{plain}
\begin{center}
  {\Large \textbf{Acknowledgment}}\\[1cm]
\end{center}

\noindent We express our deepest gratitude to our supervisor,
\textbf{Muhammad Umer Haroon}, for his continuous guidance,
encouragement, and valuable feedback throughout the course of this project.
His expertise and constructive suggestions have been instrumental in the
successful completion of our work.\\[0.5cm]

\noindent We are also sincerely thankful to our co-supervisor, \textbf{Shahzeb
Khan}, for his support, timely advice, and motivation, which helped us stay on
track and improve our project outcomes.\\[0.5cm]

\noindent We extend our appreciation to the administration and staff of
\textbf{FAST Peshawar} for their cooperation in providing domain knowledge and
feedback that shaped the system's design.\\[0.5cm]

\noindent Lastly, we would like to express our heartfelt thanks to our
families, friends, and peers for their moral support and encouragement during
this journey.\\[1.2cm]

\begin{flushright}
  Uzair Ahmad\\
  Arsalan Mateen\\
  Muhammad Sohaib
\end{flushright}

% Abstract
\newpage
\thispagestyle{plain}
\begin{center}
  {\Large \textbf{Abstract}}\\[1cm]
\end{center}

\noindent During university admission seasons, FAST-NUCES Peshawar's administrative departments receive a high volume of queries regarding NU-Test, admissions, hostels and scholarships, leading to delays and long waiting times. This report presents \textit{Mahir on Call}, a real-time voice AI web application that allows prospective students and parents to obtain instant, accurate answers directly from the university website. The system leverages Automatic Speech Recognition, Large Language Models, and
Text-to-Speech synthesis to deliver a natural conversational experience,
supplemented by live transcription for accessibility. A moderator interface
enables administrative staff to resolve queries that exceed the system's
knowledge. Evaluation demonstrates that the system substantially reduces
response latency while improving user satisfaction and accessibility for
Urdu-speaking audiences.

% Executive Summary
\newpage
\thispagestyle{plain}
\begin{center}
  {\Large \textbf{Executive Summary}}\\[1cm]
\end{center}

\noindent \textit{Mahir on Call: Advanced Elucidation Support in Local
Language} is a Final Year Project developed by students of the Department of
Computer Science, FAST-NUCES Peshawar, for the academic session 2022--2026.
The project addresses a critical operational bottleneck experienced by
university admission offices: the inability to handle hundreds of simultaneous voice and email queries during the admissions season in a timely manner.

The system is a browser-based voice AI application that a visitor can access
directly from the university website. Upon visiting the page, the user can
initiate a voice conversation in Urdu with the AI assistant. The assistant
listens, transcribes, and comprehends the spoken query using a pipeline of
Automatic Speech Recognition, Natural Language Understanding, and
a domain-specific Retrieval-Augmented Generation (RAG) engine. A
Text-to-Speech module then vocalises the generated answer in Urdu, while
simultaneously displaying a live transcript on screen.

Key features include real-time Urdu voice interaction, live transcription, and a knowledge base built from official FAST Peshawar admission documents. 
% The system is built on a modern technology stack: a React front-end, a Node.js/FastAPI back-end, OpenAI Whisper for ASR, GPT-4 for language understanding and generation, and Azure Neural TTS for speech synthesis.

The project is scoped specifically to FAST Peshawar admission queries and
demonstrates the viability of deploying conversational AI in regional
languages within resource-constrained academic environments.

% Table of Contents
\newpage
\tableofcontents

% Body pages — arabic numerals 
\newpage
\pagenumbering{arabic}
\setcounter{page}{1}

\newpage
\section{Introduction}

Conversational AI systems have rapidly moved from research laboratories into
everyday products, yet their adoption in regional-language, domain-specific
settings remains limited. Pakistan's higher-education sector faces a
particularly acute challenge during admission seasons: students and parents
speaking primarily Urdu flood university helplines with questions that are
repetitive, time-sensitive, and highly similar in nature. Staff capacity is
finite, and the resulting delays can lead poor applicant experiences.

Mahir on Call (meaning Expert on Call) is a real-time voice AI web application built specifically for the National University of Computer and Emerging Sciences, FAST. The system allows any visitor to the university website to speak a question in Urdu and receive an instant spoken answer without sending an email or navigating a complex IVR system.

A live transcript is displayed alongside the audio conversation to serve users who are deaf, hard of hearing, or who simply prefer to read.

\begin{figure}[H]
  \centering
  \includegraphics[width=0.55\textwidth]{architecture.png}
  \caption{High-level pipeline architecture of Mahir on Call}
  \label{fig:architecture}
\end{figure}

\newpage
\section{Research on Existing Products}

A survey of existing conversational AI products for higher-education and
customer-support contexts was conducted to identify gaps and opportunities.

\subsection{General Purpose Voice Assistants}

Products such as Amazon Alexa, Google Assistant, and Apple Siri are designed
for broad consumer use. While they support multiple languages, they lack
deep domain knowledge about specific institutions and cannot be customised to
answer university-specific queries without substantial engineering investment. They also do not natively support Urdu in a manner tuned for
Pakistani phonology and vocabulary.

\subsection{University Chatbot Platforms}

Several universities globally have deployed text-based chatbots (e.g., Georgia Tech's Jill Watson, built on IBM Watson). These systems answer
frequently asked questions from a curated knowledge base. However, they are
predominantly English-language, text-only systems.

\subsection{Customer-Service Voice AI}

Commercial platforms such as Dialogflow CX (Google), Amazon Lex, and Microsoft Azure Bot Service provide voice-enabled chatbot infrastructure. While powerful, these platforms require significant configuration, carry per-minute billing costs that are prohibitive for academic budgets, and do not include pre-built Urdu ASR fine-tuned on Pakistani accents.


% 3. PROJECT VISION
\newpage
\section{Project Vision}

\subsection{Problem Statement}

During the annual admissions season, the administrative offices of FAST-NUCES
Peshawar receive hundreds of emails and phone calls per day from prospective
students and their parents. The queries span topics including eligibility
criteria, the NU-Test, application procedures, fee structures, scholarship policies, and hostel availability. The volume of queries far exceeds staff capacity, leading to response delays and frustration among applicants. A significant portion of the applicant pool is Urdu speaking and is not fully comfortable with English-language digital interfaces, further narrowing their access tovtimely information.

\subsection{Business Opportunity}

A 24/7 Urdu-language voice AI service embedded in the university website would eliminate the dependency on human operators for routine, repetitive queries. This reduces operational costs, improves applicant satisfaction, and positions FAST Peshawar as a technology-forward institution. The same architecture can be extended to other FAST campuses and, in the longer term, to other Pakistani universities facing identical challenges.

\subsection{Objectives}

\begin{itemize}
  \item Develop a real-time voice AI web application accessible from any
        modern browser without additional software installation.
  \item Support Urdu language voice input and output to serve the majority of
        Pakistan's population.
  \item Provide instant, accurate answers to admission related queries using
        a RAG based knowledge base derived from official FAST Peshawar documents.
  \item Display a live transcript of the conversation for accessibility.
  \item Achieve an average response latency of under five seconds for typical
        queries.
\end{itemize}

\newpage
\subsection{Constraints}

\begin{itemize}
  % \item Urdu ASR and TTS models with
  %       Pakistani accent support are limited. The project relies on
  %       OpenAI Whisper and Azure Neural Voice, which provide acceptable
  %       but not perfect Urdu performance.
  \item The AI can only answer queries within
        its knowledge base. Queries outside scope require moderator escalation.
  \item Real-time voice processing requires a
        stable broadband connection, performance degrades on slow networks.
  \item Third-party API usage (OpenAI, Azure) introduces
        per-call costs, the system is designed to minimise token usage.
  \item  The Web Speech API is not uniformly
        supported across all browsers, Chrome and Edge offer the best
        experience.
\end{itemize}


\newpage
\section{Software Requirement Specifications}

\subsection{List of Features}

\begin{itemize}
  \item \textbf{Real-time Voice Conversation}: 
        Users can speak to the AI
        assistant and receive spoken responses in Urdu.
  \item \textbf{Live Transcription}:
        The full conversation is transcribed and displayed on screen in real time.
  \item \textbf{Urdu Language Support}:
        Operates in Urdu with English fallback.
  \item \textbf{Knowledge Base }: 
        Answers are generated from a curated RAG knowledge base of official FAST Peshawar documents.
  \item \textbf{Graceful Fallback}:
        When the AI cannot answer confidently, users had can connect with the moderator.
\end{itemize}

\subsection{Functional Requirements}

\begin{longtable}{p{2cm} p{4cm} p{8.5cm}}
  \toprule
  \textbf{ID} & \textbf{Feature}         & \textbf{Requirement} \\
  \midrule
  \endhead
  FR-01 & Voice Input           & The system shall capture microphone audio from the browser and stream it to the ASR service. \\[4pt]
  FR-02 & Speech-to-Text        & The system shall transcribe speech to text with a WER of 25\% or below on a reference test set. \\[4pt]
  FR-03 & Query Understanding   & The system shall classify the intent of the transcribed query and retrieve relevant context from the knowledge base. \\[4pt]
  FR-04 & Response Generation   & The system shall generate a contextually accurate response using the LLM. \\[4pt]
  FR-05 & Text-to-Speech        & The system shall convert the generated text response into natural Urdu audio and play it back in the browser. \\[4pt]
  FR-06 & Live Transcript       & The system shall display the ongoing conversation as text in real time. \\[4pt]
  FR-07 & Moderator Access      &  \\[4pt]
  \bottomrule
\end{longtable}

\newpage
\subsection{Non-Functional Requirements}

\begin{longtable}{p{2cm} p{3.5cm} p{9cm}}
  \toprule
  \textbf{ID} & \textbf{Category}    & \textbf{Requirement} \\
  \midrule
  \endhead
  NFR-01 & Performance          & Response latency shall be under five seconds for 90\% of queries. \\[4pt]
  NFR-02 & Availability         & The system shall target 99\% uptime. \\[4pt]
  NFR-03 & Scalability          & The system shall handle at least 20 concurrent voice sessions without degradation. \\[4pt]
  NFR-04 & Security             & All data in transit shall be encrypted using TLS 1.2 or above. No personally identifiable information shall be stored beyond the session. \\[4pt]
  NFR-05 & Usability            & A user with no technical background shall be able to initiate a voice conversation within 30 seconds of visiting the page. \\[4pt]
  NFR-06 & Accessibility        & Live transcription shall comply with WCAG 2.1 Level AA contrast and text-size guidelines. \\[4pt]
  \bottomrule
\end{longtable}


\newpage
\section{Iteration Plan}

The project was developed using an iterative approach with three major
iterations, each building on the previous and adding incremental functionality.

\subsection{Iteration 1: Voice Pipeline}

Establish the voice processing pipeline in n8n, validate the selected technologies for each stage, and deliver an English language prototype capable of answering general queries.

\subsubsection{Use Case Diagram}

Figure~\ref{fig:usecase} illustrates the primary use cases of the system
involving two actors: the User (prospective student or parent) and the
Moderator (admissions staff).

\begin{figure}[H]
  \centering
  \includegraphics[width=0.62\textwidth]{use-case.png}
  \caption{Use case diagram for Mahir on Call}
  \label{fig:usecase}
\end{figure}

\subsubsection{Use Case Descriptions}

\begin{itemize}
  \item \textbf{Start Voice Conversation}:
        The user clicks the microphone button on the web page. The browser requests microphone permission, and upon grant, begins streaming audio to the server.
  \item \textbf{View Live Transcription}:
        During and after each exchange, the user sees the conversation transcript displayed in text on screen.
  \item \textbf{End Conversation}: 
        The user clicks the stop button. The session is closed, and a full transcript is available.
  \item \textbf{Get Support}: 
        When the AI cannot answer, the system escalates to the Moderator.
\end{itemize}

% \subsubsection{Technology Stack}

% \begin{table}[H]
%   \centering
%   \caption{Technologies}
%   \begin{tabular}{p{3cm} p{9cm}}
%     \toprule
%     \textbf{Component}    & \textbf{Justification} \\
%     \midrule
%     Frontend             & React.js component-based architecture simplifies real-time state management for audio and transcript \\
%     ASR                   & OpenAI Whisper strong multilingual support including Urdu, open weights allow local deployment \\
%     LLM                   & OpenAI GPT-4 via API best-in-class instruction-following for constrained Q\&A \\
%     TTS                   & Azure Cognitive Services Neural TTS, natural prosody; low latency streaming \\
%     Backend              & Node.js and Express, non-blocking I/O suits real-time streaming \\
%     \bottomrule
%   \end{tabular}
% \end{table}


% \newpage
% \subsection{Iteration 2: LLM and Knowledge Base}

% \subsubsection{Knowledge Base Construction}

% Build and integrate a domain-specific RAG knowledge base from official FAST Peshawar admission documents including NU-Test syllabus, undergraduate admissions policy and merit criteria, scholarship programme descriptions and eligibility criteria, hostel accommodation guidelines. and campus FAQ document.

% Documents were chunked into 512-token segments, embedded using the \texttt{text-embedding-3-small} model, and stored in a Pinecone vector database. At query time, the top-$k$ ($k=5$) most semantically similar chunks are retrieved and injected into the GPT-4 prompt as context.

% \subsubsection{Urdu ASR Evaluation}

% Whisper was evaluated on 50 spoken queries recorded by native Urdu
% speakers at FAST Peshawar. The achieved Word Error Rate was 22.4\%, meeting the NFR-02 target. Common errors involved domain-specific acronyms.


\newpage
\subsection{Iteration 3: Refinement}

Integrates live transcription functionality into the user-interface to enhance accessibility and user experience. Comprehensive system and usability testing are conducted to ensure reliability and performance. Additionally, bugs are identified and resolved while continuously optimizing overall system latency.


\subsubsection{Live Transcription}

WebSocket streaming is used to push each recognised word from the ASR engine
to the browser as it is transcribed. This provides a typing-cursor like
experience for the transcript panel, giving users immediate visual feedback
that the system has heard them.

\subsubsection{Testing}

\begin{table}[H]
  \centering
  \begin{tabular}{p{5cm} p{3cm} p{5cm}}
    \toprule
    \textbf{Test Type}      & \textbf{Result} & \textbf{Notes} \\
    \midrule
    Latency      & 4.2 s avg.      & Meets NFR-01 target of $<$ 5s \\
    Query accuracy (eval set) & 82\%          & \\
    Concurrent sessions     & 20 sessions     & No degradation observed \\
    Moderator escalation    &    &  \\
    \bottomrule
  \end{tabular}
\end{table}


\newpage
\section{User Manual}

\subsection{For Prospective Students and Parents}

\subsubsection{Open NU website}
Navigate to the FAST Peshawar website, or visit the direct URL provided by the admissions office. The application will load in your browser.

\subsubsection{Allow microphone access}
When prompted by the browser, click \textbf{Allow} to grant microphone
permission. This is required for voice input.

\subsubsection{Start the conversation}
Click the microphone button. A pulsing animation indicates that the system is listening.

\subsubsection{Speak your question in Urdu}
Ask your question clearly.\\ 
\textbf{FAST Peshawar mein BS ke liye apply kaise karna hai?}

\subsubsection{Receive the answer}
The system will display your question and the AI's response as text in the
transcript panel, and will simultaneously read the answer aloud in Urdu.

\subsubsection{Continue or End conversation}
You may ask further questions. When finished, click the same microphone button to close the session.


\newpage
\section{Conclusions and Recommendations}

\subsection{Conclusions}

Mahir on Call successfully demonstrates that a real-time Urdu voice
AI system can be built and deployed for domain-specific use within an academic context. The system meets all primary objectives: it supports natural Urdu voice conversation, delivers accurate answers from a curated knowledge base, provides live transcription, and includes a human escalation mechanism.

% The project validates the use of Whisper for Urdu ASR, GPT-4 with RAG for
% domain-constrained Q\&A, and Azure Neural TTS for natural Urdu speech
% synthesis as a viable production stack.

\subsection{Recommendations}

\begin{itemize}
  \item \textbf{Knowledge base expansion:} 
        The knowledge base should be reviewed and updated annually before each admissions season to reflect
        policy changes.
  \item \textbf{Multi-campus rollout:} The modular architecture supports
        extension to other FAST campuses by swapping the knowledge base.
  \item \textbf{Analytics dashboard:} Adding analytics on query topics would
        help the admissions office identify the most common applicant concerns and improve official communication materials.
\end{itemize}


\newpage
\section*{References}
\addcontentsline{toc}{section}{References}

\begin{flushleft}

Brown, T., Mann, B., Ryder, N., Subbiah, M., Kaplan, J., Dhariwal, P., 
Neelakantan, A., Shyam, P., Sastry, G., Askell, A., \& others. (2020). 
Language models are few-shot learners. \textit{Advances in Neural Information Processing Systems}, 33, 1877--1901.

\vspace{0.5em}

Lewis, P., Perez, E., Piktus, A., Petroni, F., Karpukhin, V., Goyal, N., 
K{\"u}ttler, H., Lewis, M., Yih, W.-t., Rockt{\"a}schel, T., \& others. (2020). 
Retrieval-augmented generation for knowledge-intensive NLP tasks. 
\textit{Advances in Neural Information Processing Systems}, 33, 9459--9474.

\vspace{0.5em}

OpenAI. (2023). \textit{Whisper: Robust speech recognition via large-scale weak supervision}. 
Retrieved from \url{https://openai.com/research/whisper}

\vspace{0.5em}

Radford, A., Kim, J. W., Xu, T., Brockman, G., McLeavey, C., \& Sutskever, I. (2023). 
Robust speech recognition via large-scale weak supervision. 
In \textit{Proceedings of the 40th International Conference on Machine Learning} 
(pp.\ 28492--28518). PMLR.

\vspace{0.5em}

Pinecone Systems. (2024). \textit{Pinecone: The vector database for machine learning applications}. 
Retrieved from \url{https://www.pinecone.io}

\vspace{0.5em}

Vaswani, A., Shazeer, N., Parmar, N., Uszkoreit, J., Jones, L., Gomez, A. N., 
Kaiser, {\L}., \& Polosukhin, I. (2017). 
Attention is all you need. 
\textit{Advances in Neural Information Processing Systems}, 30.

\end{flushleft}


\newpage
\appendix
\renewcommand{\thesection}{Appendix \Alph{section}}

\section{Sample Evaluation Queries}
\label{app:queries}

The following table lists a representative sample of the 100 evaluation queries
used to measure system accuracy. Queries were collected from applicants and
admissions staff.

\begin{longtable}{p{1cm} p{8cm} p{4cm}}
  \toprule
  \textbf{No.} & \textbf{Query} & \textbf{Expected Topic} \\
  \midrule
  \endhead
  1  & FAST Peshawar mein BS admission ke liye kaun kaun se documents required hain? & Admissions documents \\
  2  & NU-Test mein kaunse subjects include hote hain? & NU-Test syllabus \\
  3  & FAST mein scholarship ke liye minimum merit ya percentage kitni honi chahiye? & Scholarship criteria \\
  4  & FAST Peshawar ke hostel mein admission ka process kya hai aur room kaise allot hota hai? & Hostel accommodation \\
  5  & FAST Peshawar ki admission fees aur semester fees kab submit karni hoti hain? & Fee payment deadlines \\
  6  & FAST Peshawar ka official admission portal link kya hai aur apply kaise karein? & Application portal \\
  7  & BS Computer Science aur Software Engineering mein kya difference hai? & Degree Programs \\
  8  & NU-Test ki tayari ke liye best strategy aur resources kya hain? & NU-Test preparation \\
  9  & Kya FAST Peshawar mein female students ke liye hostel facility available hai? & Hostel for females \\
  10 &FAST Peshawar ki merit list kab announce hoti hai aur kaise check karein? & Merit list dates \\
  \bottomrule
\end{longtable}


\end{document}
